\section{A denominator for each short section at every finite place}

% BEGIN R76V1
\paragraph{R76V1. The valuation formula on the prescribed frame.}
Fix a single short generator on the identity cover with data
$(\Lambda,N,P,x_0)$, and use the finite-place test $(V,\iota,e,y)$ of
(ICL.1). There is a finite extension of complete DVRs $V\subset V'$,
of ramification index $b$, on which the actual short value has
(ICL.13) with $\zeta=1$. The integral chart there can be taken
with compatible jet frames and with the torsion uniformization of
its semiabelian model. These are the two requirements of
$\mathrm{Good3}$.

Here are the choices in the field extension. The given integral frame
supplies a semiabelian model with the indicated polarized generic
fiber, integral eigenvectors and theta generator. The torus of its
special fiber splits after a finite separable residue extension:
its character module is a finite-rank lattice with continuous
Galois action through a finite group. Lift that residue extension
to the complete unramified trait. Use Mumford--Raynaud
uniformization for the resulting split semiabelian scheme with
its polarized degeneration data. For the associated Raynaud
extension, cubical line, and periods, see K.-W.~Lan,
\emph{Arithmetic Compactifications of PEL-Type Shimura Varieties}
(Princeton, 2013), Theorems~4.2.1.14 and 4.4.16. Its
one-motive description gives the torsion map: a $k$-torsion
point is a lift in the Raynaud extension whose $k$-multiple
belongs to the period lattice. Two lifts give the same point
exactly when they differ by a period. This gives the torsion
comparison over an algebraic closure of the fraction field. To
pass to any algebraically closed extension, take the algebraic
closure of the base field inside that extension. Modulo periods
there are only finitely many lift classes: for each residue of a
period modulo $k$, its division points form a torsor under the
$k$-torsion of the Raynaud extension. These torsors and the
torsion on the abelian variety are finite \'etale in characteristic
zero. Their geometric points and the identifications between them
therefore persist under the larger extension. The group law and
endomorphisms are induced before quotienting (Lan, \S5.1.1).
For the pairing, Lan's Theorem~5.2.3.14 and \S5.2.6 identify the
pairings on these lifts by computing the commutator of the two
theta translations for the polarized biextension. This is the
Weil pairing of the quotient with the same order convention.
The symmetric rigidification and
the endomorphism $w_0$ are transported in this construction.
This is the PEL torsion description in $\mathrm{Good3}$. It
applies to the base change of the model of $y$. Equivalently,
any integral model used for its transported chart has the same
semiabelian identity component, by uniqueness of the semiabelian
model with the given generic fiber. For the chart through the
identity section, smoothness supplies formal coordinates on the
completion of the group, and the rigidified biextension has
compatible generators on the jets: lift a generator over each
successive nilpotent thickening and normalize its axes as in
CNC5. Its group, action and line tensors are integral. The
comparison to the marked frame retains the original integral
Lie directions and section generator under base change.

Then split the finite \'etale $N^2$-torsion of the generic abelian
variety by a finite field extension. Integral closure gives a
complete DVR with finite residue field. Base change preserves the
split uniformization, its torsion description, and compatible
jet frames. The framed-point comparison throughout is the base
change of $y$ with its eigenvectors and its actual generator;
in particular the evaluation is $\iota'(\mathrm{ev}_I(y,g))$,
where $\iota'$ is the extension of fraction fields. The valuations
satisfy $\mathrm{ord}_{V'}(\iota'a)=b\,\mathrm{ord}_V(a)$,
so normalization by $be$ retains the original valuation.

All $N^2$-torsion is now rational ($\mathrm{AllRat}$). Since
$M\mid N$, the $N$-division points of the marked $M$-torsion
lie in this group. Project to the prescribed half-primary part
using the coprime level ideals. Enumerate the isotropic subgroup
of the split $N^2$-torsion cut out by $N\Lambda^\vee$: the
native lattice computation of T73C1 gives $[\Lambda:NS]^r$
tuples, where $S=\mathcal O_E^n$. Their geometric injectivity
persists under field extension. The symmetric odd-order lifts
of the theta translations supply the sections $F_{qib}$ of
$\mathrm{IsUUOver}$. They are defined over this field by
uniqueness, or by comparing with the same construction over its
algebraic closure. Evaluating A76T1--A76T2 on these data gives
(ICL.13), including its factors $\beta_H$, $\mathrm{mIdx}$,
$|\mathcal Q|$ and $N^{\deg P}$. The symbol $\mathrm{sqN}$
merely identifies the two integer exponents $(N:\mathbb Z)^2$
and $(N^2:\mathbb N)$ of the same tensor power.
% END R76V1

% BEGIN R76V2
\paragraph{R76V2. The translated columns have integral germs.}
Put $k=N^2$. Each $z_{ib}$ and each half-division point $t_{qb}$
is killed by $k$. The latter follows from $Nt_{qb}$ being a
marked $M$-torsion point and $M\mid N$. The section
$\mathrm{sqN}(F_{qib})$ is obtained from the integral generator
by the rigidified multiplication pullback along $[N]$ to power
$N^2$, then the two symmetric translations by $z_{ib}$ and
$t_{qb}$ at that power. These are exactly the three operations
defining reachability of this section, including when $N=1$.

We recall why reachability gives regularity even in the presence
of a toric part. Use the integral Fourier expansion for the
rigidified relatively ample theta line on the formal semiabelian
model (Lan, \S4.3.1, preceding Remark~4.3.1.1 and
equation~(4.3.1.7)). Modulo each positive power of the uniformizer,
sections decompose into the direct sum of the character spaces
of sections on the abelian part. The decompositions commute with
reduction, and their inverse limit is the completed direct sum.
They consequently detect vanishing on each thickening. The
associated period relation specifies the translation of the
character lines. We apply these statements to the specific
multiplication pullback and odd symmetric lifts as follows.
In the split degeneration expansion, let $Q$ be
the valuation of the period trivialization on its character
lattice. The weight-$k$ condition at lower bound $a$ is that
the coefficient with index $\chi$ have valuation at least
$a+Q(\chi)/k$, in its specified integral line. These inequalities
are valid after a valued algebraic closure and can be tested
there. The integral theta frame of weight one gives $a=0$:
periodicity reduces to the coefficient at zero, which is integral
with the prescribed rigidification. A normalized multiplication
pullback changes the index from $\chi$ to $N\chi$, takes $k$ to
$kN^2$, and gives zero off the image indices. The correction
$Q(N\chi)/(kN^2)=Q(\chi)/k$ and the rigidified integral line
comparison preserve the bound. The symmetric torsion operator
reindexes by an affine translation. In the inequality with
$Q(\chi)/k$, its index-linear change cancels the polarization
pairing of the translation; the remaining shift is a constant
independent of $\chi$. Indeed this is the polarization identity
for the quadratic term and the character of the period in the
translation law. An odd symmetric lift has finite order. Iteration
returns to the same family of coefficients, so the shift is zero
for a nonzero expansion with finite lower endpoint. More generally
the set of allowed bounds is downward closed and periodic with
period the order times the shift. If that period is nonzero and
the set contains zero it is all of $\mathbb Q$, which also proves
the claim. Any root of unity in comparing the two lifts has
valuation zero.

Thus the normalized Fourier coefficients of each column meet the
integral extension criterion on the semiabelian model. In detail,
choose a common scalar clearing a generic section on this
quasi-compact model. Its integral expansion coefficients are
divisible by the scalar's uniformizer power because their
unscaled coefficients satisfy the valuation criterion (the
polarization quadratic form is nonnegative). Reduction order by
order in the integral Fourier expansion then makes the entire
section divisible by this power: if all coefficients vanish
modulo that power, so does the section on the formal degeneration.
Removing the scalar gives a regular section with its original
generic restriction. Pull this section back through every jet and
divide by the compatible integral generator there. Its unique
coefficients commute with reduction to lower orders, yielding an
integral germ for each prescribed column on the chart of R76V1.
% END R76V2

% BEGIN R76V3
\paragraph{R76V3. One factor per derivative on any integral chart.}
Write $D_E=\operatorname{disc}(E)$ and $a=2D_E$. For an integral
chart with integral column germs and integral directions, a
homogeneous harmonic polynomial $P$ of degree $d$ with coefficients
in the valuation ring satisfies
\[
             a^d\,\mathrm{hv}(F_1,\ldots,F_n;P)\in V'.
\]
This estimate applies in all residue characteristics. The
exponent $d$ concerns the order of a jet of the section and is
independent of the ramification index of $V'$.

To prove the estimate, start with a canonical gauge $h$ for the
integral formal biextension, over the fraction field, and write $B$
for its mixed bilinear coefficient at zero. Let $A(t)$ be the
integral tangent action of $t\in\mathcal O_E$. Comparing mixed
coefficients in the symmetry and adjoint cocycles shows that
$B^t-B$ and $A(t)^tB-BA(\bar t)$ have integral entries. The
mixed coefficient of the logarithm of a normalized unit vanishing
to first order here is the mixed coefficient of the unit itself.
Choose an integral basis and its trace-dual basis. The latter is
integral after multiplication by $D_E$, by the adjugate formula.
Write their resulting tensor as $\sum x_i\otimes y_i$. Its
product is $D_E$, and $(t\otimes1)\sum x_i\otimes y_i=
(1\otimes t)\sum x_i\otimes y_i$ by the trace pairing. Hence
\[
   G=\sum_i A(x_i)^t B A(\bar y_i),\qquad
   B_0=(G+G^t)/(2D_E)
\]
satisfies $A(t)^t B_0=B_0 A(\bar t)$, $B_0^t=B_0$, and
$a(B-B_0)$ is integral. For the last assertion subtract $G$
from $D_EB$ termwise using the integral adjoint defects, then
use the integral symmetry defect. Subtract the symmetric balanced
form $B_0$ in formal logarithmic coordinates from $h$. This
preserves all canonical-gauge equations and replaces its mixed
coefficient by $B-B_0$. Denote the new gauge by $h'$.

Set $q(x)=h'(x,[w_0]x)$ and let $D_u$ be the formal invariant
derivative with integral tangent vector $u$. One has $q(0)=0$,
$D_u q(0)=0$, and $aD_uD_vq$ is an integral series. For the
last assertion, the derivatives of both addition defects of $h'$
are integral: each is the logarithmic derivative of a normalized
integral unit, hence integral by division by that unit. Expand
$h'(x+_F y,[w_0](x+_F y))$ by the two partial addition laws.
Differentiate once at $y=0$ in direction $v$. Modulo integral
series the result is the sum of the two axis derivatives of $h'$,
paired with $v$ and $A(w_0)v$. Each axis derivative is additive
modulo terms with integral derivatives by the opposite partial
addition law. Its invariant derivative therefore agrees modulo
integral series with its derivative at the origin. Consequently
$D_uD_v q$ agrees modulo integral series with the constant
$u^t(B-B_0)A(w_0)v+v^t(B-B_0)A(w_0)u$. Multiplication by $a$
gives the assertion.

In the subring generated by the integral series and all
$aD_uq$, every constant coefficient is integral, and each
$aD_u$ preserves this subring by the preceding second-derivative
bound. Starting with an integral germ $f$, differentiate
$f\exp(kq)$, scaling by $a$ at each step. After removing the
common exponential the new factor at each step is
$aD_u g+k(aD_u q)g$, still in this subring. The exponential
has constant coefficient one. A jet of order $j$ thus has
integral constant coefficient after multiplication by $a^j$.
The sum of the column orders of every monomial of $P$ is $d$.
Multiplying the column evaluations and summing with the integral
coefficients gives the displayed bound in the gauge $h'$. The
harmonic contraction is independent of the canonical gauge
and chart by G76J1--G76J3, so it is the original value of the
given connection.
% END R76V3

% BEGIN R76V4
\paragraph{R76V4. Global clearing for every original valuation test.}
Choose $h_P\in\mathbb Z_{>0}$ with
$h_PP\in\mathcal O_K[\text{letters}]$. There are finitely
many coefficients in the fixed number field $K$, so one integer
works in all its embeddings. For this section put $d=\deg P$ and
\[
 s_g=\beta_H(\Lambda)\,[\Lambda\cap S:NS]^r
       |\mathcal Q|\,(2|D_E|N)^d h_P \quad\in K^\times .
\]
If $P=0$, the section is zero and any nonzero scalar works.
Otherwise its row degrees make it homogeneous of degree $d$.
The identity-cover weights are roots of unity, integral under
every embedding in the tests. Apply R76V3 to the columns of
R76V2 and the polynomial $h_PP$. In R76V1 the displayed $s_g$
cancels the entire scalar denominator of the short value.
Each of the finitely many remaining signed terms, including
its Fourier weight, belongs to $V'$. Thus
$\mathrm{ord}_{V'}(\iota'\mathrm{ev}_I(y,s_gg))\geq0$.
Division by the positive ramification index gives the same sign
over the original $V$. The scalar was chosen before the place,
test, embedding, ramification index and frame.

This proves existence of the denominator required for the
full-lattice implication of ICL19 for each fixed generator.
In (ICL.14) the explicit lattice index and nearly all factors
of $N$ must cancel within each primary torsion sum. The present
proof clears them in $s_g$, and hence makes no assertion of
that sharper estimate for an individual $X_q$.
% END R76V4
